\subsection{空中计算与通信聚合融合}\label{subsec:aircomp}

数字正交上传需要服务器分别恢复每个客户端的更新，因而通常为终端分配不同的时隙或子载波，避免数字码字相互干扰。服务器完成解码后再进行加权，而FedAvg等模型聚合本质上就是对客户端更新求加权和。无线多址信道又具有信号叠加特性：若终端把更新映射为模拟符号并同时发送，接收波形便包含这些符号的和。空中计算（over-the-air computation，AirComp）据此把“分别传输后求和”改为“传输过程中求和”\cite{Amiri2019OTA,Zhu2020BroadbandAnalog,ZhuLi2020IntegratedAirComp}。两种机制如图~\ref{fig:digital-aircomp}所示。

\begin{figure*}[t]
  \centering
  \resizebox{0.94\textwidth}{!}{\begin{tikzpicture}[
  font=\small,
  >=Latex,
  panel/.style={draw=black!55, rounded corners=2pt, minimum width=7.9cm, minimum height=3.55cm},
  client/.style={draw, rounded corners=1.5pt, minimum width=1.45cm, minimum height=0.52cm, align=center, fill=black!3},
  block/.style={draw, rounded corners=1.5pt, minimum width=1.65cm, minimum height=0.58cm, align=center, fill=black!7},
  sum/.style={draw, circle, minimum size=0.7cm, inner sep=0pt, fill=black!7},
  flow/.style={->, line width=0.55pt}
]
  \node[panel] (leftpanel) at (-4.15,0) {};
  \node[panel] (rightpanel) at (4.15,0) {};
  \node[font=\bfseries] at (-4.15,1.43) {数字正交传输};
  \node[font=\bfseries] at (4.15,1.43) {空中模拟聚合};

  \node[client] (ld1) at (-7.15,0.72) {终端 1};
  \node[client] (ld2) at (-7.15,0) {终端 2};
  \node[client] (ld3) at (-7.15,-0.72) {终端 \(K\)};
  \node[block] (dec1) at (-4.55,0.72) {逐用户解码};
  \node[block] (dec2) at (-4.55,0) {逐用户解码};
  \node[block] (dec3) at (-4.55,-0.72) {逐用户解码};
  \node[sum] (lsum) at (-1.55,0) {\(\sum\)};
  \draw[flow] (ld1) -- node[above, font=\scriptsize] {资源 1} (dec1);
  \draw[flow] (ld2) -- node[above, font=\scriptsize] {资源 2} (dec2);
  \draw[flow] (ld3) -- node[above, font=\scriptsize] {资源 \(K\)} (dec3);
  \draw[flow] (dec1) -- (lsum);
  \draw[flow] (dec2) -- (lsum);
  \draw[flow] (dec3) -- (lsum);
  \draw[flow] (lsum) -- ++(0.72,0) node[right, font=\scriptsize] {聚合值};
  \node[font=\scriptsize, align=center] at (-4.15,-1.34) {独立时隙或子载波：先可靠解码，再在应用层求和};

  \node[client] (ra1) at (1.12,0.72) {终端 1};
  \node[client] (ra2) at (1.12,0) {终端 2};
  \node[client] (ra3) at (1.12,-0.72) {终端 \(K\)};
  \node[sum] (channel) at (4.18,0) {\(\sum\)};
  \node[block] (scale) at (6.25,0) {接收缩放\(a_t\)};
  \draw[flow] (ra1) -- node[above, sloped, font=\scriptsize] {\(h_{1,t}\beta_{1,t}x_{1,t}\)} (channel);
  \draw[flow] (ra2) -- node[above, font=\scriptsize] {\(h_{2,t}\beta_{2,t}x_{2,t}\)} (channel);
  \draw[flow] (ra3) -- node[below, sloped, font=\scriptsize] {\(h_{K,t}\beta_{K,t}x_{K,t}\)} (channel);
  \draw[flow] (channel) -- node[above, font=\scriptsize] {\(y_t\)} (scale);
  \draw[flow] (scale) -- ++(0.82,0) node[right, font=\scriptsize] {\(\widehat{s}_t\)};
  \draw[flow] (4.18,1.0) -- node[right, font=\scriptsize] {噪声 \(n_t\)} (channel);
  \node[font=\scriptsize, align=center] at (4.15,-1.34) {同频同时发送：同步、信道估计与功率对齐后在信号域求和};
\end{tikzpicture}
}
  \caption{数字正交上传与空中聚合的机理差异\protect\cite{Amiri2019OTA,Zhu2020BroadbandAnalog,ZhuLi2020IntegratedAirComp}}
  \label{fig:digital-aircomp}
\end{figure*}

为说明信道如何影响加权和，下面采用经过相位补偿的等效实数信道模型。对第 \(t\) 轮某一模型坐标，参与集合为 \(\mathcal S_t\)，终端 \(k\) 发送归一化更新 \(x_{k,t}\)，目标聚合权重为 \(q_{k,t}\)，则
\begin{subequations}\label{eq:aircomp}
  \begin{align}
    y_t &= \sum_{k\in\mathcal S_t}h_{k,t}\beta_{k,t}x_{k,t}+n_t,\label{eq:aircomp-receive}\\
    \widehat{s}_t &= a_t y_t,\qquad
    s_t=\sum_{k\in\mathcal S_t}q_{k,t}x_{k,t},\label{eq:aircomp-target}\\
    e_t &= \widehat{s}_t-s_t
    =\sum_{k\in\mathcal S_t}
      \bigl(a_t h_{k,t}\beta_{k,t}-q_{k,t}\bigr)x_{k,t}+a_t n_t.\label{eq:aircomp-error}
  \end{align}
\end{subequations}
式~\eqref{eq:aircomp-receive}表示接收端看到所有终端信号与噪声的叠加。信道系数 \(h_{k,t}\) 描述幅度变化，终端预编码 \(\beta_{k,t}\) 用于发送前补偿，\(n_t\) 是接收噪声。式~\eqref{eq:aircomp-target}中，接收缩放 \(a_t\) 调整整体幅度，\(\widehat{s}_t\) 是实际估计，\(s_t\) 是目标加权和。式~\eqref{eq:aircomp-error}进一步说明：聚合误差由系数未对齐和缩放后的噪声组成，理想对齐要求 \(a_t h_{k,t}\beta_{k,t}\approx q_{k,t}\)。这正是后文功率控制、设备选择和接收机设计需要减小的误差。

数字方案需要分离多用户信号，AirComp则把受控叠加直接用作计算结果。它的收益来自无线资源复用，而不是无损压缩更新比特。

对齐并非总能实现。信道反演令预编码系数与 \(h_{k,t}\) 近似成反比，但深衰落终端可能需要超过其最大发射功率。此时只能截断预编码、降低公共幅度或暂停该终端。信道估计不准会使 \(a_t h_{k,t}\beta_{k,t}\) 偏离目标权重；载波频率、相位或符号时刻不同步，也会破坏同一坐标上的叠加。接收噪声则直接进入式~\eqref{eq:aircomp-error}。这些因素形成的是聚合误差，不能简单依靠应用层重传消除。

AirComp适合终端多、聚合函数可加、允许小幅模拟误差，且基站能获得较准信道信息并维持同步的场景。宽带及时变信道研究通过功率控制、设备选择和接收缩放降低聚合均方误差；稀疏一比特和分层AirComp还结合压缩或网络组织，但会叠加量化、信道与层级误差\cite{Tegin2023TimeVaryingOTA,Oh2024SparseOneBitOTA,Azimi2024HierarchicalOTA}。

若深衰落频繁、同步难以保证，或任务要求可靠恢复每个客户端更新，数字上传的独立解码与差错控制更稳妥。AirComp主要减少多用户接入时间，不能理解为无条件、零误差或零能耗的聚合。

\subsection{分层联邦学习与网络拓扑优化}\label{subsec:hierarchical-fl}

终端直接连接中心云时，模型更新要跨越接入网和回传网。远距离链路增加时延、流量和掉线风险，云端还要等待分散终端。分层联邦学习把客户端按小区、位置或边缘关联划分为簇，由邻近边缘服务器先聚合。多数高频更新因而停留在端—边短链路，只有簇模型以较低频率经过边—云回传\cite{Liu2020HierarchicalFL,Abad2020HierarchicalCellular}。“端—边高频、边—云低频”表示一个云端轮次内进行多次簇内聚合，随后才同步云端，如图~\ref{fig:hierarchical-aggregation}所示。

\begin{figure*}[t]
  \centering
  \resizebox{0.62\textwidth}{!}{\begin{tikzpicture}[
  font=\small,
  >=Latex,
  cloud/.style={draw, rounded corners=2pt, minimum width=2.15cm, minimum height=0.62cm, align=center, fill=black!10},
  edge/.style={draw, rounded corners=1.5pt, minimum width=1.85cm, minimum height=0.58cm, align=center, fill=black!6},
  client/.style={draw, rounded corners=1.2pt, minimum width=1.25cm, minimum height=0.5cm, align=center, fill=black!2},
  cluster/.style={draw=black!55, dashed, rounded corners=2pt, inner sep=5pt},
  flow/.style={->, line width=0.55pt}
]
  \node[cloud] (cloud) at (0,3.15) {云端服务器\\全局聚合};
  \node[edge] (e1) at (-2.25,1.65) {边缘服务器 1};
  \node[edge] (e2) at (2.25,1.65) {边缘服务器 \(M\)};
  \node[client] (c11) at (-3.1,0.25) {终端 1};
  \node[client] (c12) at (-1.4,0.25) {终端 \(K_1\)};
  \node[client] (c21) at (1.4,0.25) {终端 1};
  \node[client] (c22) at (3.1,0.25) {终端 \(K_M\)};

  \node[cluster, fit=(e1)(c11)(c12), label={[font=\scriptsize]below:边缘簇 1}] (g1) {};
  \node[cluster, fit=(e2)(c21)(c22), label={[font=\scriptsize]below:边缘簇 \(M\)}] (g2) {};

  \draw[flow] (c11) -- (e1);
  \draw[flow] (c12) -- (e1);
  \draw[flow] (c21) -- (e2);
  \draw[flow] (c22) -- (e2);
  \draw[flow] (e1) -- node[above left, font=\scriptsize, align=center] {边—云同步\\每 \(H\) 轮一次} (cloud);
  \draw[flow] (e2) -- node[above right, font=\scriptsize, align=center] {边—云同步\\每 \(H\) 轮一次} (cloud);
  \node[font=\scriptsize, align=center] at (0,0.92) {端—边交换：高频本地更新与簇内聚合};
\end{tikzpicture}
}
  \caption{端—边—云双时间尺度聚合\protect\cite{Liu2020HierarchicalFL,Abad2020HierarchicalCellular,Liu2023HierarchicalQuantization}}
  \label{fig:hierarchical-aggregation}
\end{figure*}

为区分三种时间尺度，令 \(t\) 表示云端轮次，\(r=1,\ldots,H\) 表示其中的边缘聚合序号，\(e=0,\ldots,E_{k,m,t}^{(r)}\) 表示本地SGD步数。这里 \(H\) 只表示每个云端轮次包含的边缘聚合次数，不表示本地步数。设系统有 \(M\) 个边缘簇，第 \(m\) 个簇的客户端集合为 \(\mathcal C_{m,t}\)，本地轨迹为 \(\bm u_{k,m,t}^{(r,e)}\)，最终本地模型记为 \(\bar{\bm u}_{k,m,t}^{(r)}=\bm u_{k,m,t}^{(r,E_{k,m,t}^{(r)})}\)。两级聚合可写为
\begin{subequations}\label{eq:hierarchical-aggregation}
  \begin{align}
    \bm w_{m,t}^{(0)}&=\bm w_t,\qquad
    \bm u_{k,m,t}^{(r,0)}=\bm w_{m,t}^{(r-1)},\label{eq:hierarchical-init}\\
    N_{m,t}&=\sum_{k\in\mathcal C_{m,t}}n_k,\qquad
    \bm w_{m,t}^{(r)}
    =\sum_{k\in\mathcal C_{m,t}}\frac{n_k}{N_{m,t}}
      \bar{\bm u}_{k,m,t}^{(r)},\label{eq:edge-aggregation}\\
    N_t&=\sum_{m=1}^{M}N_{m,t},\qquad
    \bm w_{t+1}
    =\sum_{m=1}^{M}\frac{N_{m,t}}{N_t}\bm w_{m,t}^{(H)}.\label{eq:cloud-aggregation}
  \end{align}
\end{subequations}
式~\eqref{eq:hierarchical-init}表示云模型先下发到边缘，并成为客户端的初始模型。式~\eqref{eq:edge-aggregation}表示边缘服务器按样本数汇总簇内模型。式~\eqref{eq:cloud-aggregation}表示完成 \(H\) 次边缘聚合后，云端再把各簇模型合成为下一轮全局模型。三式依次对应图~\ref{fig:hierarchical-aggregation}中的云端下发、端—边训练和边—云同步。

增大 \(H\) 会在两次云端同步之间执行更多簇内聚合，因此减少边—云同步频率。各簇也会沿自身数据分布独立演化更久；簇间数据差异较大时，边缘模型可能向不同方向移动，形成簇间漂移，并需要更多云端轮次重新对齐。客户端分簇会影响数据代表性，边缘负载不均会产生新的慢节点，终端移动还会改变集合 \(\mathcal C_{m,t}\) 并增加模型迁移与控制通信。

分层方法适合多小区蜂窝、工业物联网和云—边—端集群，其收益取决于回传瓶颈、簇内数据分布和边缘可用性。量化分层可同时选择聚合间隔和消息精度；异构物联网方案还要处理异步到达、边缘负载和客户端选择\cite{Liu2023HierarchicalQuantization,Tang2023AdaptiveGradient,Wang2023HierarchicalIoT}。端—边链路也可结合压缩或AirComp，并用资源分配平衡各簇\cite{Azimi2024HierarchicalOTA}。组合收益不能简单相加，需要分别核算两级比特、时延与能耗，并报告压缩误差、信道失真和层级漂移。
